\documentclass[uplatex,dvipdfmx]{jlreq}
%\documentclass{jlreq}
\usepackage{amsmath,amssymb, listings, graphicx, float, hyperref, subcaption}

\lstset{
  frame=single,
  basicstyle=\ttfamily,
}

\begin{document}

\title{システム提案書}

\author{KRYMAS}

\date{\today}

\maketitle

\setcounter{tocdepth}{3}
\tableofcontents
\clearpage


\subsection{ログイン・プロフィール}
        \begin{table}[h]
        \centering
        \caption{初回利用設定機能}
        \begin{tabular}{c|c|c|c|c|c}
        \hline
        変数名 & データ型 & 長さ & NOT NULL & その他制約 & 説明 \\
        \hline
         password & VARCHAR & 20 & 〇 &  &パスワード \\
         petName& VARCHAR & 10 & 〇 &  &ペット名\\
         mailAddress & VARCHAR & 50 & 〇 &  &メールアドレス\\
        \hline
        \end{tabular}
        \end{table}

        \begin{table}[h]
        \centering
        \caption{ログイン・ログアウト}
        \begin{tabular}{c|c|c|c|c|c}
        \hline
        変数名 & データ型 & 長さ & NOT NULL & その他制約 & 説明 \\
        \hline
         ID & VARCHAR & 50 & 〇 &  &ログインID \\
         password& VARCHAR & 20 & 〇 &  &パスワード\\
        \hline
        \end{tabular}
        \end{table}

%%ここからやる
        \begin{table}[h]
        \centering
        \caption{パスワード管理機能}
        \begin{tabular}{c|c|c|c|c|c}
        \hline
        変数名 & データ型 & 長さ & NOT NULL & その他制約 & 説明 \\
        \hline
         password & VARCHAR & 20 & 〇 &  &パスワード \\
         mailAddress & VARCHAR & 50 & 〇 &  &メールアドレス\\
         auCode& VARCHAR & 4 & 〇 &  &認証コード\\
        \hline
        \end{tabular}
        \end{table}

        \begin{table}[h]
        \centering
        \caption{プロフィール管理機能}
        \begin{tabular}{c|c|c|c|c|c}
        \hline
        変数名 & データ型 & 長さ & NOT NULL & その他制約 & 説明 \\
        \hline
         length & INTERGER & 3 & 〇 & 0より大きい &身長 \\
         weight & INTERGER & 3 & 〇 & 0より大きい &体重\\
         gender& NCHAR & 3 & 〇 & 選択肢を制限 &性別\\
         experience& VARCHAR & 255 & 〇 & 選択肢を制限 &運動経験\\
         likefood& VARCHAR & 255 &  &  &食事の好み\\
         request& VARCHAR & 255 &  &  &要望（アレルギー等）\\
         freetime& VARCHAR & 255 &  &  &空き時間\\
         petName& VARCHAR & 10 & 〇 &  &ペット名\\
         target& VARCHAR & 255 & 〇 &  &目標設定\\
        \hline
        \end{tabular}
        \end{table}

\end{document}
